% Generated draft: requires mathematical and source review.
\documentclass[11pt]{article}
\usepackage[margin=25mm]{geometry}
\usepackage{fontspec}
\setmainfont{Latin Modern Roman}
\usepackage{amsmath,amssymb,mathtools,unicode-math}
\setmathfont{Latin Modern Math}
\usepackage{enumitem,needspace,xurl,hyperref}
\setlength{\parindent}{0pt}
\setlength{\parskip}{5pt}
\setlength{\emergencystretch}{3em}
\newcommand{\N}{\mathbb{N}}
\newcommand{\Z}{\mathbb{Z}}
\newcommand{\Q}{\mathbb{Q}}
\newcommand{\R}{\mathbb{R}}
\newcommand{\C}{\mathbb{C}}
\newcommand{\F}{\mathbb{F}}
\newcommand{\A}{\mathbb{A}}
\newcommand{\PP}{\mathbb{P}}
\newcommand{\E}{\mathbb{E}}
\newcommand{\eps}{\varepsilon}
\newcommand{\dd}{\,\mathrm{d}}
\begin{document}
\section*{UnsolvedMath ID: 30003614}
\textbf{Source problem code:} catalogue entry
\subsection{Definitions and mathematical statement}
Let \(K\) be a field finitely generated over \(\Q\), fix an algebraic closure \(\overline K\), and write \(\Gamma_K=\operatorname{Gal}(\overline K/K)\). Let \(X\) be a smooth projective variety over \(K\), and let \(\ell\) be a prime. Interpret the source's notation \(H^i(X)(n)\) as the continuous \(\ell\)-adic Galois representation
\[
V_{i,n,\ell}(X)=H^i_{\mathrm{\acute et}}\bigl(X\times_K\overline K,\Q_\ell\bigr)\otimes_{\Q_\ell}\Q_\ell(n).
\]
Here \(\Q_\ell(1)\) is the one-dimensional representation defined by the \(\ell\)-adic cyclotomic character, \(\Q_\ell(n)\) is its \(n\)-th tensor power for \(n>0\), \(\Q_\ell(0)\) is the trivial representation, and \(\Q_\ell(n)\) for \(n<0\) is the dual of \(\Q_\ell(-n)\). A finite-dimensional \(\Q_\ell\)-representation of \(\Gamma_K\) is semisimple if every \(\Gamma_K\)-invariant \(\Q_\ell\)-linear subspace has a \(\Gamma_K\)-invariant complementary subspace; equivalently, it is a direct sum of irreducible representations.

\textbf{Question.} For every such \(K\) and \(X\), every integer \(i\geq 0\), every integer \(n\), and every prime \(\ell\), is \(V_{i,n,\ell}(X)\) semisimple as a representation of \(\Gamma_K\)?

\textbf{Interpretive qualification.} The source does not explicitly specify the cohomology theory, coefficient field, or ranges of \(i,n\). The formulation above uses the customary interpretation as rational \(\ell\)-adic geometric étale cohomology, with nonnegative cohomological degree and integral Tate twist. These interpretations require human confirmation.

\textbf{Supplied status qualification.} The registry labels this problem partially solved. Its dated assessment reports that Moonen proved in characteristic zero that the Tate cycle-class conjecture implies semisimplicity, while the general assertion above remains unresolved. This is supplied curation, not an independently verified literature assessment; the precise scope and hypotheses of that reported implication are not supplied.
\subsection{Short English statement}
Does every rational \(\ell\)-adic cohomology group of a smooth projective variety over a finitely generated field of characteristic zero, with any integral Tate twist, decompose into irreducible Galois representations? The cohomology notation and index ranges use the customary interpretation of the supplied statement. The partially solved label and reported conditional result are supplied, dated assessments requiring review.
\subsection{Sources}
\begin{itemize}
\item Supplied local UnsolvedMath registry record: permanent ID \(30003614\), title \emph{Galois Semisimplicity of Cohomology over Finitely Generated Fields}, source problem code \texttt{catalogue entry}, in the Oberwolfach Reports Open Problems collection. The record gives proposed year 2017 and names no individual proposer. Its labels \texttt{corrected\_verified} and \texttt{partially\_solved} are supplied registry classifications, not independent verification here.
\item Supplied source citation: \emph{Algebraic Geometry: Birational Classification, Derived Categories, and Moduli Spaces} (2018), Oberwolfach Reports, pp. 2703--2767, DOI 10.4171/owr/2017/45. \url{https://doi.org/10.4171/owr/2017/45}. The background describes Oberwolfach Reports as documenting mathematical workshops and their open-problem discussions.
\item Supplied literature reference: Moonen, \emph{A remark on the Tate conjecture}. \url{https://arxiv.org/abs/1709.04489}.
\item Supplied literature reference: \emph{Oberwolfach Report 45 2017}. \url{https://doi.org/10.4171/owr/2017/45}.
\item Literature checked at: \texttt{2026-08-21}. This exact date belongs to the supplied dated curation, described in the background as literature triage derived from the supplied OWR curation. Its assessment reports a characteristic-zero implication from the Tate cycle-class conjecture to semisimplicity and describes the general assertion as unresolved. Neither the sources nor the time-sensitive status and cited claims have been independently verified here.
\end{itemize}
\end{document}
